\section{Scope, baseline, and the new results}
\label{sec:scope}

The linked repository contains a consolidated manuscript, not merely
the earlier collection of experimental identities. At the pinned
revision it has already incorporated proofs of the Gaussian weight-six
double identities, the displayed weight-four triples, the mixed $S_4$
relation, several rank formulas, and the cubic and quartic
trilogarithm ladders. It also contains sharp real-order kernel theorems,
integral distribution modules, reflected moments, and a substantial
zero-geometry programme. Those advances must be treated as the
starting point. Repeating them would not answer the request for a
continuation.

We audited the active questions against the source chapters and the
later reports preserved beside them. Four directions admit substantive
progress with proofs that can be stated precisely and checked
independently. They share a method: first find the exact object that
controls the question, and then separate its analytic or algebraic
structure from numerical diagnostics. For Euler errors that object is
a two-point kernel difference. For depth it is the osculating span of
binary forms. For radial transitions it is a convergent implicit
equation, together with certified low-order derivatives. For reflected
moments it is a saddle integral with an oscillatory amplitude and a
signed cut representation.

\subsection{A theorem and status ledger}

\begin{longtable}{@{}>{\raggedright\arraybackslash}p{0.18\textwidth}>{\raggedright\arraybackslash}p{0.35\textwidth}>{\raggedright\arraybackslash}p{0.40\textwidth}@{}}
\toprule
Direction & Status in the audited material & Contribution of this article\\
\midrule
\endhead
Euler errors & Order-dependent bounds; sharp critical constant;
axis magnitude maximum; unrestricted remainder comparison open.
& A sharp all-order remainder theorem, its optimal global constant,
an exact rational budget, positive moment identities, and a moving-order
phase law with an asymptotically optimal constant and an explicit
lower asymptotic for its approach.\\[5pt]
Formal depth & Exact calculations, including a weight-six obstruction
and special transformations.
& An all-weight formula for any number of projective directions;
six-map versus full-linear comparison; exact certificates for known
Nielsen reductions, with literature attribution.\\[5pt]
Radial geometry & Both signs of the quartic coefficient on the quadratic
threshold; a diagnostic simultaneous zero.
& Exact isolation, negative sixth coefficient, a quarter-power law,
a two-extremum local unfolding, and a fully explicit rational example.\\[5pt]
Reflected coefficients & Fixed reflected exponent asymptotics and
fixed-exponent optimal truncation.
& A uniform $m=O(\sqrt{k})$ transition, Gaussian damping, explicit
first correction, a growing-exponent half-term law, and a small proportional saddle theorem.\\[5pt]
Special-value candidates & $S_6$ and the revised $S_8$ candidate remain
conjectural despite certified small residuals.
& Their evaluation status is unchanged; the new formal obstruction
specifies a limitation of a particular reduction method.\\
\bottomrule
\end{longtable}
\addtocounter{table}{-1}

The labels ``new'' and ``open'' in this article refer to the audited
repository state. The general tools---Euler transformation, gamma
integrals, free-Lie character theory, Hermite interpolation, analytic
implicit inversion, and the saddle method---are classical. We do not
claim that a broad literature search could establish priority for every
consequence. The article supplies proofs of the stated refinements
and records the primary sources used for their background.

\subsection{Conventions and the interpretation of identities}

Multiple polylogarithms use decreasing nested indices:
\[
 \Li_{s_1,\ldots,s_d}(z_1,\ldots,z_d)
 =\sum_{n_1>\cdots>n_d\geq1}
        \frac{z_1^{n_1}\cdots z_d^{n_d}}
             {n_1^{s_1}\cdots n_d^{s_d}}
\]
where the series converges. A one-variable specialization means
$\Li_{s_1,\ldots,s_d}(z,1,\ldots,1)$. Real-order $F_{a,b}$ is
defined explicitly in Section~\ref{sec:euler}, so no convention for
nonintegral-depth words is needed. All analytic continuations use the
principal branch unless a simply connected path domain is specified.
Statements at $i$ concern analytic or Abel values where the ordinary
Fourier series fails to converge.

Formal identities in a shuffle quotient have a different meaning.
Their vanishing proves an exact congruence modulo the stated product
and endpoint-constant relations. A nonzero class proves the absence
of a certificate using those relations. It does not prove numerical
independence of periods and does not exclude other functional
equations. Section~\ref{gorbit:sec:main} gives the exact comparison space
before drawing any depth conclusion.

For parameter estimates, terminating decimal endpoints are rational
numbers. They can therefore be part of a finite exact proof when
every operation is outward rounded and every transcendental input
has a proved remainder bound. By contrast, a printed approximate
root or a many-digit agreement from quadrature is diagnostic unless
an enclosure theorem is also supplied.

\subsection{What is reproducible}

The archive includes this source and its compiled PDF, independent
exact verifiers for the finite claims, numerical checks that compare
different representations, and the resulting machine-readable data.
The new statements are proved in the text; the code has three
different roles, explicitly separated in the verification section:
finite arithmetic proof certificates, finite corroboration of a
general algebraic theorem, and nonrigorous asymptotic diagnostics.

The repository was read at the pinned revision without applying
changes to it. Suggested chapter insertions and wording replacements
are supplied separately. The package can thus be reviewed and
integrated against a concrete baseline without silently altering
historical reports or their recorded status.
